\chapter{Simulación en LTspice}

\section{Polarización en directa}
\begin{figure}[H]
    \centering
    \includegraphics[width=0.6\textwidth]{img/CircuitoDiodoPolDir.png}
    \caption{Circuito de diodo 1N4007 polarizado en directa}
    \label{fig:CircuitoDiodoPolDir}
\end{figure}
\begin{tcolorbox}[
colback=gray!10,
colframe=black,
width=\textwidth,
boxrule=0.5pt,
title={Modelo SPICE del diodo 1N60},
arc=0pt
]
\begin{verbatim}
.dc V1 -2 10 100m
.model 1N4007 D (IS=7.02767n RS=0.0341512 N=1.80803 EG=1.05743 XTI=5 
BV=1000 IBV=5e-08 CJO=1e-11 VJ=0.7 M=0.5 FC=0.5 TT=1e-07 mfg=OnSemi 
type=silicon)
\end{verbatim}

\end{tcolorbox}
\begin{tcolorbox}[
colback=gray!10,
colframe=black,
width=\textwidth,
boxrule=0.5pt,
title={Modelo SPICE del diodo 1N60},
arc=0pt
]
\begin{verbatim}
.dc V1 -2 10 100m
.model 1N60 D (RS=50.0559 EG=1.11 XTI=3 VJ=0.8 M=0.1 FC=0.5 
BV=40 CJO=1.75p IS=1.37773u N=2 TT=100n type=germanium)
\end{verbatim}

\end{tcolorbox}

\subsection{Diodo 1N4007 polarizado en directa}
\begin{figure}[H]
\centering
\begin{tikzpicture}
\begin{axis}[
    width=14cm,
    height=9cm,
    ymin=-0.1,
    xmax=1,
    xmin=-0.1,
    xlabel={$V_D$ [V]},
    ylabel={$I_D$ [mA]},
    grid=major,
    minor tick num=1,
    grid=both,
]
\addplot[
    ultra thick,
    blue
]
table[
    col sep=tab,
    header=true,
    x index=1,
    y index=2,
    y expr=\thisrowno{2}*1000
] {dataTex/1N4007PolDir};


\end{axis}
\end{tikzpicture}
\caption{Curva del diodo 1N4007 en directa, obtenida por simulación}
\end{figure}

\subsection{Diodo 1N60 polarizado en directa}
\begin{figure}[H]
\centering
\begin{tikzpicture}
\begin{axis}[
    width=14cm,
    height=9cm,
    ymin=-0.1,
    xmax=1,
    xmin=-0.1,
    xlabel={$V_D$ [V]},
    ylabel={$I_D$ [mA]},
    grid=major,
    minor tick num=1,
    grid=both,
]
\addplot[
    ultra thick,
    red
]
table[
    col sep=tab,
    header=true,
    x index=1,
    y index=2,
    y expr=\thisrowno{2}*1000
] {dataTex/1N60PolDir};

\end{axis}
\end{tikzpicture}
\caption{Curva del diodo 1N60 en directa, obtenida por simulación}
\end{figure}

\subsection{Comparaciones}
\begin{figure}[H]
\centering
\begin{tikzpicture}
\begin{axis}[
    width=15cm,
    height=10cm,
    ymin=-0.01,
    xmax=1,
    xmin=-0.1,
    xtick distance=0.1,
    xlabel={$V_D$ [V]},
    ylabel={$I_D$ [mA]},
    grid=both,
    minor tick num=1,
    legend pos=north west
]

% Primer diodo
\addplot[
    ultra thick,
    blue
]
table[
    col sep=tab,
    header=true,
    x index=1,
    y index=2,
    y expr=\thisrowno{2}*1000
] {dataTex/1N4007PolDir};

% Segundo diodo
\addplot[
    ultra thick,
    red
]
table[
    col sep=tab,
    header=true,
    x index=1,
    y index=2,
    y expr=\thisrowno{2}*1000
] {dataTex/1N60PolDir};

\legend{1N4007, 1N60}

\end{axis}
\end{tikzpicture}
\caption{Comparación de curvas }
\end{figure}

Se observó que ambos diodos comienzan a conducir a partir de cierto voltaje umbral, pero con comportamientos diferentes. El diodo 1N60 presenta una conducción más gradual y comienza a conducir con tensiones menores, mientras que el 1N4007 permanece casi sin 
corriente hasta aproximadamente 0,6–0,7 V, donde la corriente aumenta bruscamente. Esto se debe a que el 1N60 es un diodo de germanio y el 1N4007 es de silicio.
También se observa que la curva del 1N60 no sube tan recta como la del 1N4007. Esto significa que, al aumentar el voltaje, la corriente del 1N60 aumenta más lentamente, mientras que en el 1N4007 la corriente crece muy rápido después de alcanzar su voltaje de conducción.

\section{Polarización en inversa}
Para este caso se usaran los mismos modelos SPICE de Polarización directa. 
\begin{tcolorbox}[
colback=gray!10,
colframe=black,
width=\textwidth,
boxrule=0.5pt,
arc=0pt
]
\begin{verbatim}
    Figura 3.5:
        .dc V1 0 -1200 100m
    Figura 3.9: 
        .dc V1 10 -60 100m
\end{verbatim}

\end{tcolorbox}

\subsection{Diodo 1N4007}
\begin{figure}[H]
\centering
\begin{tikzpicture}
\begin{axis}[
    width=14cm,
    height=9cm,
    ymin=-5,
    xmax=20,
    xmin=-1200,
    xtick distance=200,
    xlabel={$V_D$ [V]},
    ylabel={$I_D$ [mA]},
    grid=major,
    minor tick num=1,
    grid=both,
]
\addplot[
    thick,
    blue
]
table[
    col sep=tab,
    header=true,
    x index=1,
    y index=2,
    y expr=\thisrowno{2}*1000
] {dataTex/1N4007PolInv};
\end{axis}
\end{tikzpicture}
\caption{1N4007 en inversa, obtenido por simulación}
\end{figure}

\subsection{Diodo 1N60}

\begin{figure}[H]
\centering
\begin{tikzpicture}
\begin{axis}[
    width=14cm,
    height=9cm,
    ymin=-5,
    xmax=10,
    xmin=-60,
    xtick distance=10,
    xlabel={$V_D$ [V]},
    ylabel={$I_D$ [mA]},
    grid=major,
    minor tick num=1,
    grid=both,
]
\addplot[
    thick,
    red
]
table[
    col sep=tab,
    header=true,
    x index=1,
    y index=2,
    y expr=\thisrowno{2}*1000
] {dataTex/1N60PolInv};
\end{axis}
\end{tikzpicture}
\caption{1N60 en Inversa, obtenido por simulación}
\end{figure}





\begin{figure}[H]
\centering
\begin{tikzpicture}
\begin{axis}[
    width=14cm,
    height=9cm,
    ymin=-5,
    xmax=10,
    xmin=-1050,
    xlabel={$V_D$ [V]},
    ylabel={$I_D$ [mA]},
    grid=major,
    minor tick num=1,
    grid=both,
]

% Primer diodo
\addplot[
    ultra thick,
    blue
]
table[
    col sep=tab,
    header=true,
    x index=1,
    y index=2,
    y expr=\thisrowno{2}*1000
] {dataTex/1N4007PolInv};

% Segundo diodo
\addplot[
    ultra thick,
    red
]
table[
    col sep=tab,
    header=true,
    x index=1,
    y index=2,
    y expr=\thisrowno{2}*1000
] {dataTex/1N60PolInv};

\legend{1N4007, 1N60}

\end{axis}
\end{tikzpicture}
\caption{Comparación 1N4007 vs 1N60}
\end{figure}

Mediante la simulación en polarización inversa se logró observar la diferencia en la tensión de ruptura de ambos diodos, obteniéndose aproximadamente \(BV=40\,V\) para el 1N60 y \(BV=1000\,V\) para el 1N4007, valores correspondientes a los 
parámetros proporcionados por el fabricante en los modelos utilizados para la simulación. Esto permitió comprender el comportamiento de los diodos en la zona negativa a altas tensiones.

\section{Comportamiento del diodo en función de la temperatura}
Para este caso se usaran los mismos modelos SPICE de Polarización directa.
\begin{tcolorbox}[
colback=gray!10,
colframe=black,
width=\textwidth,
boxrule=0.5pt,
arc=0pt
]
\begin{verbatim}
    Figura 3.8:
        .step temp list 20 100 125 
        .dc V1 0 3 0.01
    Figura 3.9:
        .step temp list 20 100 125 
        .dc V1 -3 1 0.01
\end{verbatim}

\end{tcolorbox}

\begin{tcolorbox}[
colback=gray!10,
colframe=black,
width=\textwidth,
boxrule=0.5pt,
arc=0pt
]
\begin{verbatim}
    Figura 3.10: 
        .step temp list 20 100 125
        .dc V1 -2 3 0.01 
\end{verbatim}

\end{tcolorbox}

\subsection{Temperatura en el diodo 1N4007}

\begin{figure}[H]
\centering
\begin{tikzpicture}
\begin{axis}[
    width=14cm,
    height=9cm,
    ymin=-0.2,
    xmax=3,
    xmin=0,
    xlabel={$V_1$ [V]},
    ylabel={$I_D$ [mA]},
    grid=major,
    minor tick num=1,
    grid=both,
]

% 20°C
\addplot[blue, ultra thick]
table[
col sep=tab,
x index=0,
comment chars={S},
y expr=\thisrowno{1}*1000
] {dataTex/1N4007Temp-dir-20};

% 100°C
\addplot[green!70!black, ultra thick]
table[
col sep=tab,
x index=0,
comment chars={S},
y expr=\thisrowno{1}*1000
] {dataTex/1N4007Temp-dir-100};

% 125°C
\addplot[red, ultra thick]
table[
col sep=tab,
x index=0,
comment chars={S},
y expr=\thisrowno{1}*1000
] {dataTex/1N4007Temp-dir-125};

\legend{20°C,100°C,125°C}

\end{axis}
\end{tikzpicture}
\caption{Curvas para diferentes temperaturas en directa(1N4007)}
\end{figure}


\begin{figure}[H]
\centering
\begin{tikzpicture}
\begin{axis}[
    width=14cm,
    height=9cm,
    ymin=-4,
    ymax=4,
    xmax=1,
    xmin=-3,
    enlarge y limits=false,
    xlabel={$V_1$ [V]},
    ylabel={$I_D$ [muA]},
    grid=major,
    minor tick num=1,
    grid=both,
    scaled y ticks=false,
    yticklabel style={
        /pgf/number format/fixed,
        },
]

% 20°C
\addplot[blue, ultra thick]
table[
col sep=tab,
x index=0,
y expr=\thisrowno{1}*1000000
] {dataTex/1N4007Temp-inv-20};

% 100°C
\addplot[green!70!black, ultra thick]
table[
col sep=tab,
x index=0,
y expr=\thisrowno{1}*1000000
] {dataTex/1N4007Temp-inv-100};

% 125°C
\addplot[red, ultra thick]
table[
col sep=tab,
x index=0,
y expr=\thisrowno{1}*1000000
] {dataTex/1N4007Temp-inv-125};

\legend{20°C,100°C,125°C}

\end{axis}
\end{tikzpicture}
\caption{Curvas para diferentes temperaturas en inversa (1N4007)}
\end{figure}

\subsection{Temperatura en el diodo 1N60}

\begin{figure}[H]
\centering
\begin{tikzpicture}
\begin{axis}[
    width=14cm,
    height=9cm,
    ymin=-1,
    xmax=3,
    xmin=-2,
    xlabel={$V_1$ [V]},
    ylabel={$I_D$ [mA]},
    grid=major,
    minor tick num=1,
    grid=both,
]

% 20°C
\addplot[blue, ultra thick]
table[
col sep=tab,
x index=0,
comment chars={S},
y expr=\thisrowno{1}*1000
] {dataTex/1N60Temp-20};

% 100°C
\addplot[green!70!black, ultra thick]
table[
col sep=tab,
x index=0,
comment chars={S},
y expr=\thisrowno{1}*1000
] {dataTex/1N60Temp-100};

% 125°C
\addplot[red, ultra thick]
table[
col sep=tab,
x index=0,
comment chars={S},
y expr=\thisrowno{1}*1000
] {dataTex/1N60Temp-125};

\legend{20°C,100°C,125°C}

\end{axis}
\end{tikzpicture}
\caption{curvas para diferentes temperaturas  (1N60)}
\end{figure}

\subsection{Análisis de resultados}

En la Figura 3.9 se observa la característica en polarización directa del diodo para
distintas temperaturas. A medida que la temperatura aumenta, la curva se desplaza hacia
la izquierda, indicando que para un mismo voltaje se obtiene mayor corriente. Esto se
debe a la reducción de la barrera de potencial y al aumento de portadores disponibles en
la unión PN.

En la Figura 3.10 se muestra el comportamiento en polarización inversa. Se observa
una corriente de fuga que aumenta significativamente con la temperatura. Este efecto
es causado por la generación térmica de portadores minoritarios, lo que incrementa la
conducción inversa del diodo.

Para el caso del diodo 1N60, se aprecia además una mayor sensibilidad a la temperatura
en comparación con el 1N4007, presentando un aumento más notorio de la corriente
inversa y una variación más visible en la zona de conducción. Esto se relaciona con las
características propias de los diodos de germanio, los cuales poseen menores tensiones
de conducción y mayores corrientes de fuga.

En conclusión, la temperatura incrementa la corriente tanto en directa como en inversa,
afectando directamente el comportamiento ideal del diodo.

\section{Circuitos recortadores con diodos zener}

\subsection{Recortador con diodos zener modelo 1}
\begin{figure}[H]
    \centering
    \includegraphics[width=0.7\textwidth]{img/recortadores-zener-1.png}
    \caption{}
    \label{fig:recortadores-zener-1}
\end{figure}

\begin{figure}[H]
\centering
\begin{tikzpicture}
\begin{axis}[
    width=14cm,
    height=9cm,
    ymin=-23,
    ymax=23,
    xmin=0,
    xmax=60,
    xlabel={$t$ [ms]},
    ylabel={$V$ [v]},
    grid=major,
    minor tick num=1,
    grid=both,
]

% V1 tension de entrada
\addplot[
    ultra thick,
    blue
]
table[
    col sep=tab,
    header=true,
    y index=1,
    x expr=\thisrowno{0}*1000
] {dataTex/recortadores1-v1};

% V0 ,tension salida
\addplot[
    ultra thick,
    red
]
table[
    col sep=tab,
    header=true,
    y index=1,
    x expr=\thisrowno{0}*1000
] {dataTex/recortadores1-v0};

\legend{V1, V0}

\end{axis}
\end{tikzpicture}
\caption{Recorte de una señal senoidal mediante diodos Zener}
\end{figure}

\subsection{Recortador con diodos zener modelo 2}
%Aqui va figura

\begin{figure}[H]
    \centering
    \includegraphics[width=0.7\textwidth]{img/recortadores-zener-2.png}
    \caption{}
    \label{fig:recortadores-zener-2}
\end{figure}

\begin{figure}[H]
\centering
\begin{tikzpicture}
\begin{axis}[
    width=14cm,
    height=9cm,
    ymin=-23,
    ymax=23,
    xmin=0,
    xmax=60,
    xlabel={$t$ [ms]},
    ylabel={$V$ [v]},
    grid=major,
    minor tick num=1,
    grid=both,
]

% V1 tension de entrada
\addplot[
    ultra thick,
    blue
]
table[
    col sep=tab,
    header=true,
    y index=1,
    x expr=\thisrowno{0}*1000
] {dataTex/recortadores2-v1};

% V0 ,tension salida
\addplot[
    ultra thick,
    red
]
table[
    col sep=tab,
    header=true,
    y index=1,
    x expr=\thisrowno{0}*1000
] {dataTex/recortadores2-v0};

\legend{V1, V0}

\end{axis}
\end{tikzpicture}
\caption{Recorte de una señal senoidal mediante diodos Zener}
\end{figure}

En las simulaciones se comprobó el funcionamiento de los diodos zener como recortadores
de tensión en corriente alterna. La señal de salida quedó limitada a valores cercanos a las
tensiones zener utilizadas, observándose además que la orientación de los diodos modifica
los niveles de recorte positivo y negativo. Los resultados obtenidos coincidieron con el
comportamiento teórico  desarrollados en la sección 2.2 .